\section{膜量子化回顾}\label{two}

\subsection{问题}\label{problem}

我们的出发点是 $2n$ 维实辛流形 $M$，辛形式为 $\omega$。假定 $M$（称为\textbf{相空间}）配有一个“\textbf{预量子线丛}”，即复线丛 $\frL\to M$，带有一个曲率为 $\omega$ 的酉联络（这里暂缓讨论“亚辛校正”，那需要对这一说法稍作修正）。这样的 $\frL$ 存在，当且仅当 $\omega$ 的周期取值于 $2\pi\Z$。若再有 $H^1(M,\UU(1))=0$，则 $\frL$ 在同构意义下唯一。量子化将依赖于 $\frL$ 的选取。

粗略地说，量子化的想法是：给这组数据配上一个 Hilbert 空间 $\H$，使 $M$ 上的函数成为 $\H$ 上的算符，函数的 Poisson 括号对应算符的换位子，即若 $\h f$ 是函数 $f$ 对应的算符，则
\be\label{goal} [\h f,\h g]=-\i\hbar \h {\{f,g\} }. \ee
（这里 $\{~,~\}$ 是函数的 Poisson 括号，$[~,~]$ 是算符的换位子，$\hbar$ 是 Planck 常数。）$M$ 上的实函数应当对应 $\H$ 上的 Hermite 算符。$\H$ 的维数应当大致等于辛体积 $\frac{1}{(2\pi)^n}\int_M e^\omega$；特别地，这一维数无限当且仅当 $M$ 的体积无限。

没有额外结构时，量子化问题没有自然的解。除函数的一些特殊类之外，不可能（在“大小”合适的 Hilbert 空间里）满足方程 \eqref{goal}。量子力学教科书称之为“\textbf{算符排序问题}”，但人们并不总强调这个问题没有一般解。从经典力学到量子力学的过渡存在反常。

经典力学的自然对称性由 $M$ 上的实函数（视为哈密顿函数）生成。一个函数 $f$ 经由向量场 $v^i=\omega^{ij}\frac{\partial f}{\partial x^j}$ 生成 $M$ 的一个辛变换（这里 $x^i$ 是 $M$ 上的局部坐标，$\omega=\frac{1}{2}\omega_{ij}\d x^i \d x^j$，$\omega^{ij}\omega_{jk}=\delta^i{}_k$）。这样的向量场称为哈密顿向量场，它们生成一个群 $\G_0$。此群的元素称为\textbf{恰当辛同构}。\footnote{若 $M$ 的第一 Betti 数正，则它还有非恰当的辛同构；其生成元是多值哈密顿函数。注意 $\G_0$ 的元素自动同伦于恒等映射，因为这样的元素粗略地说是把某个哈密顿函数的作用指数化得到的。与恒等映射不同伦的 $M$ 的辛同构，在我们把 $\G_0$ 换成 $\G$ 之后马上就会包含进来。}

从 $f$ 过渡到 $v^i=\omega^{ij}\partial_j f$ 时丢掉了一些信息，因为 $v^i$ 在给 $f$ 加常数时不变。为了尽可能精确地让经典力学与量子力学对应，应当追踪这个常数；做法是考虑并非 $M$ 本身、而是 $M$ 连同其预量子线丛 $\frL$ 的对称性。对 $(M,\frL)$ 的对称性构成一个群 $\G$，我们把它视为经典力学的对称群。

具体地说，$\G$ 的元素是一个辛同构 $\varphi:M\to M$ 连同同构 $\beta:\frL\to \varphi^*(\frL)$。把同构作为定义的一部分可能不常见，但实际上这能导致经典力学与量子力学的对称性之间尽可能接近（尽管仍很不完美）的匹配。$\G$ 中含恒等元的连通分支 $\G_1$ 是恰当辛同构群 $\G_0$ 被 $\UU(1)$ 的中心扩张：
\be\label{centext}1\to \UU(1)\to \G_1\to \G_0\to 1. \ee
这里 $\UU(1)$ 在 $M$ 上平凡地作用，而在 $\frL\to M$ 的纤维上作用为乘以模 1 的复常数 $e^{\i\alpha}$。$\G$ 一般还有与恒等元不连通的其它分支。

与之相对，量子力学的自然对称群是 Hilbert 空间 $\H$ 的酉变换群 $\U$。群 $\G$ 与 $\U$ 并不相同，而且二者的差异不只是一个中心扩张的问题。$\G$ 含有相空间 $M$ 的拓扑信息，而 $\U$ 不含。这就是量子场论意义下的反常：经典系统的对称群在量子化时被修改了。

例如，取 $M=\R^{2n}$ 及其标准辛结构。假定已给定“$\R^{2n}$ 上次数 $\leq 1$ 的多项式”这一概念的含义。重要的是，并不存在这一概念的自然含义；给出这一概念，等价于给出一组特殊的坐标 $\vec x =(x^1,x^2,\cdots, x^{2n})$，它在一个变换 $\vec x \to A\vec x +\vec b$（$A$ 为辛矩阵，$b$ 为常数）之下有定义。这样的变换称为仿射线性变换；若 $M$ 被赋予一组在仿射线性变换下唯一确定的坐标，我们就说 $M$ 具有\textbf{仿射线性结构}。若 $\R^{2n}$ 上给定了一个仿射线性结构，则任何非线性辛同构都会把这个仿射线性结构映到另一个仿射线性结构。因此 $\R^{2n}$ 上仿射线性结构的空间是无限维的。

在给定的仿射线性坐标系中，生成仿射线性变换的哈密顿函数是次数 $\leq 2$ 的多项式。对次数 $\leq 2$ 的多项式没有反常。这是一个众所周知的事实，只是可能没有人用完全这样的措辞。它等价于说：相空间坐标的仿射线性变换在量子力学中实现为对称性。\footnote{\label{cext}这里涉及一个二重覆盖。相空间的线性变换 $\vec x\to A\vec x$ 构成群 $\Sp(2n,\R)$；相应的量子对称群是它的二重覆盖。在 $\R^2$ 中，二次哈密顿函数的一个例子是谐振子哈密顿量 $H=\frac{1}{2}(p^2+q^2)$，它生成 $\R^2$ 的一个旋转。量子对称群是经典对称群的二重覆盖，因为 $H$ 的能级是半整数，于是对应于 $\G_0$ 中恒等元的 $\exp(2\pi\i H)$ 在量子 Hilbert 空间上作用为 $-1$。}（例如，Fourier 变换交换坐标与动量。）然而，对次数 $>2$ 的多项式会出现反常；任何算符排序或添加低阶项都无法使它们的量子换位子与经典 Poisson 括号一致。这个相当古老的结果 \cite{G46,V51,Todorov} 在附录 \ref{anomaly} 中对 $n=2$ 的情形给出了显式演示。

一旦给定仿射线性结构，就可以用熟悉的方式量子化 $\R^{2n}$，例如把 $x$ 分成 $p$ 与 $q$（选取使得 $\omega = \sum_i \d p_i \d q^i$），并定义由 $q$ 的 $\ltwo$ 函数构成的 Hilbert 空间。标准论证表明，如何作这种分解其实无关紧要；所得的 Hilbert 空间之间有自然的等同（差一个整体因子 $\pm 1$，见脚注 \ref{cext}）。这一陈述实际上等价于：次数 $\leq 2$ 的多项式的换位子没有反常。然而，从不同的仿射线性结构出发构造出的 Hilbert 空间之间，即使允许差 $c$-数因子，也没有自然的等价。若有，最终就会证明经典力学与量子力学的对称群 $\G$ 与 $\U$ 相同，而事实并非如此。

所以任何量子化方法都需要在 $M$ 上附加辛结构与预量子线丛之外的结构。这就是 Ludwig Faddeev 过去常说量子化是“一门艺术，而非科学”的原因。不过，这些问题可能不为人熟悉是有充分理由的。尽管有一些值得注意的例外，物理学中的相空间通常有足够的附加结构（例如自然的仿射线性结构），使得这里描述的问题并不突出。
